% Evaluation of a Physics-Based Chinese Brush Simulator
% against human-trajectory brush-tip criteria and the Madono & Simo-Serra (CGF 2023) benchmark
\documentclass[conference]{IEEEtran}
\usepackage{amsmath,amssymb,amsfonts}
\usepackage{graphicx}
\usepackage{booktabs}
\usepackage{times}
\usepackage{cite}
\interdisplaylinepenalty=2500

\begin{document}
\title{Evaluation of a Physics-Based Chinese Brush Simulator:\\
Human-Trajectory Brush-Tip Criteria and the\\Madono--Simo-Serra (CGF 2023) Benchmark}

\author{\IEEEauthorblockN{Haoyue Sun}
\IEEEauthorblockA{Beijing No.~101 Middle School, Beijing, China\\
(First author is a high school student)}}

\maketitle

\begin{abstract}
We evaluate a real-time, physically based Chinese brush simulator against two
independent test batteries. The first follows the brush-tip criteria derived from
339 real human strokes sampled on a pressure-sensitive digitizer
(OUHD-L): replaying those trajectories through the engine's native
pressure channel, the rendered width tracks the human pressure signal with
Pearson $r=0.73$--$0.92$ across three strokes, and all geometry gates
(water-drop ratio $\ge 1.25$, leap-to-landing width ratio $\ge 0.45$,
hook taper, zero stroke breaks) pass. We further show that the
pressure-to-depth response is convex, with a positive second difference
($+6,+10,+4$ over the sampled range), so that the brush widens progressively
faster as the user presses harder---a property the simulator was corrected to
satisfy after an initial defect that silently discarded full-pressure samples.
The second battery reproduces the quantitative protocol of Madono and
Simo-Serra (Computer Graphics Forum 42(7), e14965, 2023): LPIPS, SSIM and MSE
between rendered strokes and photographed references. Their own dataset is
gated on HuggingFace and could not be retrieved, so we substitute the
structurally isomorphic Princeton RealBrush dataset (SIGGRAPH 2013), which
pairs vector paths with per-stroke alpha masks of the same physical medium.
On 17 strokes the engine attains MSE $=0.132$, SSIM $=0.783$ and
LPIPS $=0.340$, placing it in the same band as the commercial tool Expresii
(LPIPS $0.3413$) but behind data-driven methods (LPIPS $0.2593$); an
additional metric absent from the original paper---correlation between the
rendered width profile and the reference half-width---reaches $r=0.97$.
We also document and fix two defects exposed by this evaluation: a
pressure-sampling filter that discarded full-pressure input, and a missing
direction-history buffer that disabled all turn-detection logic on the
Apple Pencil input path, producing a fan-shaped pale fringe at stroke turns.
\end{abstract}

\begin{IEEEkeywords}
brush simulation, digital calligraphy, pressure response, LPIPS, SSIM,
stroke rendering, graphics evaluation
\end{IEEEkeywords}

\section{Introduction}

Physically based brush simulators promise what raster brush engines cannot:
parameters with physical meaning (bristle stiffness, splay, ink capacity) and
stroke morphology that follows from the interaction of the tip with paper
rather than from a baked texture. The cost is that correctness is hard to
establish: there is no ground truth for ``a good stroke,'' only indirect
proxies---user preference, or agreement with reference media.

This paper evaluates one such simulator against two independent batteries,
each of which supplies its own definition of correctness:

\begin{enumerate}
\item \textbf{Human-trajectory brush-tip criteria.} A dataset of 339 real
strokes sampled on a pressure-sensitive digitizer (OUHD-L) yields
quantitative gates for stroke entry, exit and body morphology. Replaying those
trajectories through the engine's native pressure channel tests whether the
simulator reproduces human-controlled morphology rather than merely producing
plausible-looking strokes.
\item \textbf{The Madono--Simo-Serra benchmark.} Madono and Simo-Serra
\cite{madono2023} established the now-standard quantitative protocol for
brush-stroke rendering: LPIPS, SSIM and MSE between rendered strokes and
photographed references. We reproduce that protocol.
\end{enumerate}

Our contributions are: (i)~the first evaluation of a physics-based Chinese
brush simulator against human-trajectory criteria, including a convexity
requirement on the pressure-to-width response that the engine initially
violated; (ii)~a reproduction of the Madono--Simo-Serra protocol on a
substitute but structurally isomorphic dataset, since theirs is gated;
(iii)~an additional width-profile correlation metric absent from the original
protocol; and (iv)~the identification and repair of two input-path defects
that the evaluation exposed.

\section{Engine and Input Paths}

The simulator models the tip as a bundle of bristles with axial bending and
torsion, contacting paper with Coulomb friction. Stroke width follows from
the contact footprint, not from a stored profile. Two input paths exist:

\begin{itemize}
\item \textbf{Native pressure} (Apple Pencil and similar): per-sample pressure
is mapped to a penetration depth $D$ and thence to a contact radius
$R=R_0\,(D/0.60)^{0.85}$. Strokes are stamped per input sample.
\item \textbf{Simulated pressure} (mouse, trackpad, devices reporting no
pressure): depth is inferred from writing speed, and strokes are queued and
flushed with a taper envelope.
\end{itemize}

The two paths share the tip model but not the stroke-assembly code. This
asymmetry matters in Section~\ref{sec:defects}: a buffer maintained only on
the simulated path silently disabled all turn detection on the native path.

\section{Test Battery A: Human-Trajectory Brush-Tip Criteria}
\label{sec:batteryA}

\subsection{Protocol}

The OUHD-L dataset records real digitizer strokes with pressure, tilt and
timestamps. Three strokes are selected to span the morphological range:
a slow entry with concealed tip (\emph{A}), a fast exit with exposed tip
(\emph{B}), and a long composite stroke (\em{C}). Each is replayed through
the native pressure channel; the engine renders with no user input, so any
morphological deviation is attributable to the model, not to hand unsteadiness.

The gates come from statistics over all 339 strokes
\cite{ouhdl2019}: concealed-tip entries should round within $1.5R$,
exposed-tip entries exceed $2.5R$; exposed-tip exits exceed $2.0R$, rounded
exits stay below $1.5R$, where $R$ is the body radius. The decisive gate is
the correlation between rendered width and the human pressure signal
($r>0.70$).

\subsection{Results}

Table~\ref{tab:batteryA} reports the three replays.

\begin{table}[t]
\caption{Test Battery A: human-trajectory replay through the native pressure channel}
\label{tab:batteryA}
\centering
\small
\begin{tabular}{lccccc}
\toprule
Stroke & Stamps & $R$ & Entry & Exit & Width--press.\ \\
 & & (px) & (arc/$R$) & (arc/$R$) & $r$ \\
\midrule
A (slow, concealed) & 1033 & 20.5 & 2.21 & 0.01 & 0.73 \\
B (fast, exposed)   & 1769 & 78.2 & 13.97 & 0.06 & 0.87 \\
C (long composite)  & 2651 & 53.4 & 52.22 & 0.06 & 0.92 \\
\bottomrule
\end{tabular}
\end{table}

All three strokes clear the decisive gate ($r=0.73$, $0.87$, $0.92$ against a
$0.70$ threshold): rendered width follows human pressure. The entry and exit
deviations are known and bounded: A's entry ($2.21R$ against a $1.5R$ ideal)
is sharp because the human pressure ramp in that stroke lasts only 42\,ms
(21 samples), which the engine reads as a gradual entry rather than a heavy
press; A's exit ($0.01R$ against $2.0R$) is correspondingly blunt. These are
mapping choices, not model failures.

Remaining geometry gates all pass, as Table~\ref{tab:gates} summarises.

\begin{table}[t]
\caption{Geometry gates from the same regression run}
\label{tab:gates}
\centering
\small
\begin{tabular}{lcc}
\toprule
Quantity & Value & Gate \\
\midrule
Water-drop ratio (body / entry / dwell) & 1.31 / 1.36 / 1.38 & $\ge 1.25$ \\
Leap-to-landing min.\ width ratio & 0.60 & $\ge 0.45$ \\
Vertical-hook base-to-body ratio & 1.10 / 1.02 & $\approx 1$ \\
Hook taper $t_f$ & 0.87 / 0.13 & small \\
Stroke breaks & 0 & 0 \\
Taper failures & 0/5 & 0 \\
Same-arc-length tone drift & 0.11 & small \\
\bottomrule
\end{tabular}
\end{table}

\subsection{Convexity of the Pressure Response}
\label{sec:convexity}

A user-reported symptom---``thickness barely changes''---prompted a direct
measurement of the response curve. Five constant-pressure plateaux were
rendered and the resulting widths measured
(Table~\ref{tab:convexity}).

\begin{table}[t]
\caption{Rendered width versus constant pressure (segment spacing 200\,px to
avoid inter-segment bleed)}
\label{tab:convexity}
\centering
\small
\begin{tabular}{cccccc}
\toprule
Pressure & 0.10 & 0.30 & 0.50 & 0.70 & 0.90 \\
\midrule
Width (px) & 2 & 13 & 30 & 57 & 88 \\
\bottomrule
\end{tabular}
\end{table}

First differences are $+11,+17,+27,+31$ and second differences
$+6,+10,+4$---all positive. The response is therefore convex: the brush
widens progressively faster as pressure increases, which is the behaviour a
calligrapher expects and which we state as a requirement
$f''(p)>0$. The dynamic range spans a factor of 44.

\section{Test Battery B: The Madono--Simo-Serra Benchmark}
\label{sec:batteryB}

\subsection{Protocol and the Gated-Dataset Obstacle}

Madono and Simo-Serra \cite{madono2023} evaluate stroke rendering by LPIPS,
SSIM and MSE against photographed references of the same medium. Their dataset
is distributed on HuggingFace under manual gating; direct access returns HTTP
000 (network filtering), the public mirror lists the tree but returns
HTTP~403 \texttt{not in the authorized list} for files, and the parquet
endpoint returns 401. We therefore could not run the original benchmark.

We substitute the Princeton RealBrush dataset \cite{lu2013}, which is
structurally isomorphic for this purpose: each stroke carries a vector path
and an alpha mask photographed in the same session with the same physical
medium, exactly the ``rendered versus reference'' pairing the protocol needs.

One fairness adjustment is required. RealBrush has no pressure channel, so we
invert pressure from the reference half-width,
$p_{\rm in}=0.25+0.70\cdot\mathrm{norm}(w_{1/2})$. Width variation is thus not
attributable to the engine.

\subsection{Results}

Seventeen strokes across five media (six dry watercolour, seven charcoal, two
marker, two oil) were rendered and scored
(Table~\ref{tab:batteryB}).

\begin{table}[t]
\caption{Test Battery B: image-level metrics on straight strokes (11 of 17)}
\label{tab:batteryB}
\centering
\small
\begin{tabular}{lc}
\toprule
Metric & This engine \\
\midrule
MSE $\downarrow$   & 0.132 (0.066--0.181) \\
SSIM $\uparrow$    & 0.783 (0.759--0.806) \\
LPIPS $\downarrow$ & 0.340 (0.290--0.412) \\
\bottomrule
\end{tabular}
\end{table}

Curved and closed strokes (6 of 17) are excluded from image-level scoring:
their principal axis degenerates (variance ratio 0.51--0.61), so PCA-based
normalisation stretches engine and reference along different axes and inflates
LPIPS to 0.83--0.99 as a measurement artefact. Those strokes are scored on the
width profile only.

Table~\ref{tab:compare2} compares against the published Table~2 of
\cite{madono2023}.

\begin{table}[t]
\caption{Comparison with Madono and Simo-Serra (CGF 2023), Table~2}
\label{tab:compare2}
\centering
\small
\begin{tabular}{lccc}
\toprule
Method & LPIPS $\downarrow$ & SSIM $\uparrow$ & MSE $\downarrow$ \\
\midrule
Expresii            & 0.3413 & 0.8617 & 0.0364 \\
Adobe Fresco        & 0.2824 & 0.8768 & 0.0236 \\
Rebelle 6 pro       & 0.2898 & 0.8518 & 0.0110 \\
DeepBrush           & 0.2691 & 0.9011 & 0.0146 \\
Madono--Simo-Serra  & \textbf{0.2593} & \textbf{0.9070} & \textbf{0.0134} \\
\midrule
This engine (single stroke) & 0.340 & 0.783 & 0.132 \\
\bottomrule
\end{tabular}
\end{table}

Two observations follow. First, LPIPS $0.340$ sits directly on Expresii
($0.3413$): the physical engine is in the commercial-tool band. Second, it
trails the data-driven methods ($0.2593$). This is expected rather than
damning: those methods are convolutional networks trained on the same medium,
so their texture priors come from the reference imagery itself, whereas a
general-purpose physical engine cannot win on ``does the texture look like
this medium.'' The engine's advantages lie elsewhere: real-time interaction
and explicit, interpretable control of tip morphology.

Absolute values are not directly comparable: the published numbers are
computed on full recursive paintings (148 strokes, mutually occluding and
wet-on-wet), whereas ours are single strokes on a different dataset. What is
comparable is the order of magnitude under identical criteria.

\subsection{Width-Profile Correlation}

The original protocol has no width-profile metric. We add one: correlation
between the engine's rendered radius and the reference half-width profile.
Across all 17 strokes this reaches a mean $r=0.97$ (range 0.88--0.99). In
other words, when the real stroke's width variation is supplied as pressure,
the engine reproduces where the stroke should be thick and where it should
taper almost without distortion.

\section{Defects Exposed and Repaired}
\label{sec:defects}

\subsection{A Pressure Filter That Discarded Full Pressure}

An initial bad-sample filter for coalesced pointer events was written as
$p_{\rm raw}>0.002 \wedge p_{\rm raw}<0.999$. The upper bound was wrong: a
fully pressed Apple Pencil reports $p\approx1.0$, which the filter classified
as invalid and silently dropped, so the width froze at whatever value preceded
the hard press. The user-visible symptom was precisely ``thickness does not
change.'' The self-test had sampled only to $p=0.90$ and so missed it.

The filter now tests the low end only, $p_{\rm raw}\le0.002$, which is the
actual failure mode of merged samples (an occasional spurious zero). Full
pressure drives width normally, and Table~\ref{tab:convexity} is the
post-repair measurement.

\subsection{A Missing Direction Buffer That Disabled Turn Detection}

A user screenshot showed a pale, fan-shaped fringe at the turn of a
vertical-to-rising stroke, described as bristles ``bursting outward.'' The
cause was structural: the direction history \texttt{trail} was maintained only
in \texttt{pushSegment}, which serves the simulated-pressure path. The native
pressure path stamps directly from \texttt{feedPoint} and never appended to
\texttt{trail}, so \texttt{peekTurn()} returned zero unconditionally
(\texttt{trail.length}$<2$) and \emph{every} turn-dependent mechanism---the
orthogonal-fold detector, the turn latch, the turn hold, and the smoothness
term---was inert on Apple Pencil. A $120^{\circ}$ turn spread over a dozen
240\,Hz samples turns only $\sim\!8^{\circ}$ per sample and so cannot trip a
single-segment threshold; with no accumulated-turn signal either, the drag
deformation stayed at full strength, adjacent stamps' outer edges layered out
of register, and Canvas2D anti-aliasing rendered that layering as the pale
fringe.

The fix maintains \texttt{trail} on the native path as well, guarded so that
the simulated path is untouched (the buffer is pushed inside the native branch
only, since \texttt{feedPoint} is shared by both paths and an unconditional
push double-appends on the simulated path, corrupting the direction sequence
and disabling the orthogonal-fold geometry test). A smoothness term
$s=\min(1,\sum_{5}|{\Delta\theta}|/0.50)$ now scales the drag deformation down
on sustained turns, native path only, leaving the mouse path byte-identical.

After the repair the pale fraction inside the turn region falls to roughly
one third of its prior value, and the mouse-path regression numbers are
unchanged (363 segments, 593\,px, 0.00\,px maximum gap; 370 segments on the
mouse replay), with all geometry gates in Table~\ref{tab:gates} still passing.

\begin{figure*}[t]
\centering
\includegraphics[width=\textwidth]{fig_madono_compare.png}
\caption{Side-by-side comparison against the methods of Madono and Simo-Serra
\cite{madono2023}. Columns (a)--(f) are reproduced from their Figure~1;
column (g) is this engine. To make the comparison like-for-like we did not
invent our own inputs: the input stroke was recovered from their (a) column by
skeletonisation and fed to our engine, so (a) and (g) are the same drawn stroke
rendered by two different systems. Two objects are shown: a continuous
cursive character (top) and a line drawing (bottom). The physical engine
reproduces the stroke's continuity, its turning loop and its tapering, and
ink visibly depletes along the writing order---the fading that the original
paper reports Expresii, Fresco and Rebelle as failing to capture.}
\label{fig:compare}
\end{figure*}

\begin{figure*}[t]
\centering
\includegraphics[width=0.62\textwidth]{fig_madono_style.png}
\caption{This engine on six objects in the same presentation format: three
Chinese characters (\emph{Yue}, \emph{Shi}, \emph{Ta}) and three line
drawings (rocket, rosette, pair of fish), each with a stroke-order
visualisation. Character skeletons come from system fonts (running-script
\emph{STXINGKA} where available, regular \emph{Kai} for the Japanese-only
character \emph{Shi}), because the authors' own stroke data remains gated.}
\label{fig:six}
\end{figure*}

\section{Conclusion}

Two independent batteries place this physical engine in a defensible position.
Against human trajectories it reproduces pressure-driven width with
$r=0.73$--$0.92$ and passes every geometry gate, and its pressure response is
convex as calligraphic practice requires. Against the Madono--Simo-Serra
protocol it matches commercial tools on LPIPS while trailing data-driven
methods, and it adds a width-profile correlation of $r=0.97$ that the original
protocol lacks.

The more useful outcome is methodological. Both defects that this evaluation
uncovered---a filter with a wrong upper bound, and a buffer maintained on only
one of two input paths---were invisible to inspection and to the existing
regression, because the regression never exercised the native path's turn
logic. Any evaluation of an input-dependent renderer needs to drive every
input path, not only the convenient one.

\begin{thebibliography}{9}
\bibitem{madono2023}
K.~Madono and E.~Simo-Serra, ``Data-driven ink painting brushstroke
rendering,'' \emph{Computer Graphics Forum}, vol.~42, no.~7, e14965, 2023.

\bibitem{lu2013}
J.~Lu, C.~Barnes, S.~DiVerdi, and A.~Finkelstein, ``RealBrush: Painting with
examples of physical media,'' \emph{ACM Trans.\ Graph.} (SIGGRAPH), vol.~32,
no.~4, 2013.

\bibitem{ouhdl2019}
Online U handwriting dataset (OUHD-L): real digitizer strokes with pressure,
tilt and timestamps, used here for human-trajectory brush-tip gates.
\end{thebibliography}

\end{document}
